\subsection{Prehilbertian spaces}

DEFINITION 4. --- A prehilbertian space is a set E with the structure of a vector space over K and with a positive hermitian form. We say that E is a real (resp. complex) prehilbertian space when K is R (resp. K is C).

Examples. --- 1) The form \((\lambda, \mu) \mapsto \overline{\lambda}\mu\) defines a prehilbertian structure on K, said to be canonical. When K is considered as a prehilbertian space, we shall always mean, unless otherwise mentioned that it has this structure.

\begin{enumerate}
\def\labelenumi{\arabic{enumi})}
\setcounter{enumi}{1}
\tightlist
\item
  Let I be an interval (bounded or not) in R, and let E be the set of regulated functions (FVR, II, p.~4) defined on I with values in C, having compact support. It is clear
\end{enumerate}

that \(\mathbf{E}\) is a vector space over \(\mathbf{C}\) ; let \(f\) be the sesquilinear form \((x, y) \mapsto \int_{1}^{\infty} \overline{x(t)} y(t) dt\) ;

it is immediate that \(f\) is a positive hermitian form on \(E\) , and hence defines a prehilbertian structure on this space.

\begin{enumerate}
\def\labelenumi{\arabic{enumi})}
\setcounter{enumi}{2}
\tightlist
\item
  Let \(n \geqslant 0\) be an integer. We define a prehilbertian space structure on the space \(\mathbf{K}^n\) , by means of the hermitian form
\end{enumerate}

\[
(x, y) \mapsto \sum_ {j = 1} ^ {n} \overline {{{{x}}}} _ {j} y _ {j}
\]

(for \(x = (x_{1}, \ldots, x_{n})\) and \(y = (y_{1}, \ldots, y_{n})\) ). When \(\mathbf{K}\) is \(\mathbf{R}\) , we see that this is just the scalar product of two vectors of \(\mathbf{R}^{n}\) (GT, VI, § 2, No.~2).

* 4) Let \(\ell^2\) (or \(\ell^2(\mathbf{N})\) ) be the set of sequences \(x = (x_n)_{n \in \mathbb{N}}\) of elements of \(\mathbf{K}\) such that \(\sum_{n=0}^{\infty} |x_n|^2\) is finite. One can show that \(\ell^2\) is a vector subspace of \(\mathbf{K}^{\mathbf{N}}\) and define a pre-

hilbertian space structure on \(\ell^2\) by means of the hermitian form \((x, y) \mapsto \sum_{n=0}^{\infty} \overline{x}_n y_n\) (cf.~V, p.~18).

\begin{enumerate}
\def\labelenumi{\arabic{enumi})}
\setcounter{enumi}{4}
\tightlist
\item
  Let \(\mathbf{E}\) be a real prehilbertian space, \(f\) the corresponding symmetric bilinear form on \(\mathbf{E}\) . Let \(\mathrm{E}_{(\mathbf{C})}\) be the vector space complexification of \(\mathbf{E}\) ; we identify \(\mathbf{E}\) with a subset of \(\mathrm{E}_{(\mathbf{C})}\) by the map \(x \mapsto 1 \otimes x\) , in such a way that every element of \(\mathrm{E}_{(\mathbf{C})}\) can be written uniquely as \(x_1 + ix_2\) with \(x_1, x_2\) in \(\mathbf{E}\) . The map \(f\) extends uniquely to a hermitian form \(f_{(\mathbf{C})}\) on \(\mathrm{E}_{(\mathbf{C})}\) ; we have,
\end{enumerate}

\[
f _ {(\mathbf {C})} \left(x _ {1} + i x _ {2}, y _ {1} + i y _ {2}\right) = f \left(x _ {1}, y _ {1}\right) + f \left(x _ {2}, y _ {2}\right) + i \left(f \left(x _ {1}, y _ {2}\right) - f \left(x _ {2}, y _ {1}\right)\right).
\]

In particular, we have

\[
f _ {(\mathbf {C})} (x _ {1} + i x _ {2}, x _ {1} + i x _ {2}) = f (x _ {1}, x _ {1}) + f (x _ {2}, x _ {2}) \geqslant 0,
\]

hence \(f_{(\mathbf{C})}\) is positive. We say that \(\mathrm{E}_{(\mathbf{C})}\) , with \(f_{(\mathbf{C})}\) is the prehilbertian space complexification of \(\mathbf{E}\) .

Whenever only one prehilbertian space structure on a vector space E is under consideration, the value, for a pair \((x, y)\) of points of E, of the hermitian form which defines the said structure is denoted by \(\langle x|y\rangle_{E}\) or simply \(\langle x|y\rangle\) , if no confusion is likely to arise. This number is called the scalar product \(^{1}\) of x and y (scalar square of x if y = x). Two vectors x, y are said to be orthogonal if \(\langle x|y\rangle = 0\) . The function \(x \mapsto \|x\| = \langle x|x\rangle^{1/2}\) is a semi-norm on the vector space E (V, p.~3); a prehilbertian space is always considered with this semi-norm assigned to it (and consequently also with the corresponding topology and uniform structure).

With these notations, in a prehilbertian space E, the Cauchy-Schwarz inequality can be written as

\[
\left| \langle x | y \rangle \right| \leqslant \| x \|\cdot \| y \|.\tag{9}
\]

Consequently, the scalar product is a continuous sesquilinear form on \(\mathbf{E} \times \mathbf{E}\) (II, p.~5, prop. 4).

In order that E be Hausdorff, it is necessary and sufficient that \(x \mapsto \|x\|\) is a norm on E; in other words, that the hermitian form \((x, y) \mapsto \langle x | y \rangle\) is positive and separating; this is equivalent to saying that 0 is the only vector of E, which is orthogonal to itself.

According to general definitions (S, IV, § 1, No.~5), an isomorphism from a prehilbertian space E onto a prehilbertian space F is a bijective linear mapping u from E onto F such that

\[
\langle u (x) | u (y) \rangle = \langle x | y \rangle\tag{10}
\]

for every x and y in E. We deduce from this that \(\|u(x)\| = \|x\|\) for all \(x \in E\) , and u is evidently an isomorphism for the topological vector space structures of E and of F; if E and F are Hausdorff, u is an isometry from E onto F. Conversely, if u is a bijective linear mapping from E onto F, such that \(\|u(x)\| = \|x\|\) for all \(x \in E\) , the polarization formulas (V, p.~2) show that u is a prehilbertian space isomorphism from E onto F.

Let \(E\) be a complex prehilbertian space, and \(\langle x|y\rangle\) the scalar product in \(E\) . On the set \(E\) , we can define a second vector space structure with respect to \(C\) , taking the same law of the additive group and for the law of external composition \((\lambda, x) \mapsto \overline{\lambda}x\) (A, II, § 1, No.~13) for this vector space structure, \((x, y) \mapsto \langle y|x\rangle\) is a positive hermitian

\(^{1}\) It may happen sometimes that we write \((x|y)\) for \(\langle y|x\rangle\) . Observe that the formula (4) of V, p.~2, takes the following equivalent forms:

(4')

\[
\langle \sum_ {i} \lambda_ {i} x _ {i} | \sum_ {j} \mu_ {j} y _ {j} \rangle = \sum_ {i, j} \overline {{{{\lambda}}}} _ {i} \mu_ {j} \langle x _ {i} | y _ {j} \rangle .\tag{4"}
\]

\[
(\sum_ {i} \lambda_ {i} x _ {i} | \sum_ {j} \mu_ {j} y _ {j}) = \sum_ {i, j} \lambda_ {i} \overline {{\mu}} _ {j} (x _ {i} | y _ {j}).
\]

form. The prehilbertian space \(\overline{E}\) obtained by assigning this new vector space structure and the new hermitian form to E, is said to be conjugate to E. An isomorphism u from E onto \(\overline{E}\) is a semi-linear mapping from E onto itself (with respect to the automorphism \(\xi \mapsto \overline{\xi}\) of C) such that \(\langle u(y)|u(x)\rangle = \langle x|y\rangle\) or \(\langle u(x)|u(y)\rangle = \overline{\langle x|y\rangle}\) (for x, y in E); such a mapping is said to be a semi-automorphism of the prehilbertian space E.

If E is a prehilbertian space, M a vector subspace of E, the restriction of the scalar product \(\langle x|y\rangle\) to \(M \times M\) is a positive hermitian form on M, which then defines a prehilbertian space structure on M; we say that this structure is induced by the structure of E, or that M is a prehilbertian subspace of E.

